\newcommand{\celllist}[1]{%
  \parbox[t]{\linewidth}{\vspace{-2.3mm}#1}%
}

\begin{table}[H]
\centering
\renewcommand{\arraystretch}{1.25}
\setlength{\tabcolsep}{6pt}

\begin{tabular}{
    >{\centering\arraybackslash}p{0.07\textwidth}
    >{\RaggedRight\arraybackslash}p{0.21\textwidth}
    >{\RaggedRight\arraybackslash}p{0.29\textwidth}
    >{\RaggedRight\arraybackslash}p{0.31\textwidth}
}
\toprule
\textbf{Section} & \textbf{Approach} & \textbf{Examples in deep learning} & \textbf{Examples from physics%
} \\
\midrule
\ref{sec:reason-dynamics}
&
solvable settings
&
\celllist{
deep linear networks,\\
kernel regression,\\
multi-index models
}
&
\celllist{
harmonic oscillator,\\
hydrogen atom,\\
Ising model
}
\\[3em]

\ref{sec:reason-limits}
&
simplifying limits
&
\celllist{
lazy vs.\ rich learning,\\
width, depth $\rightarrow \infty$,\\
small initialization
}
&
\celllist{
thermodynamic limit $(n, V \rightarrow \infty)$,\\
classical limit $(\hbar \rightarrow 0)$,\\
hydrodynamic limit $(\vk, \omega \rightarrow 0)$
}
\\[3em]

\ref{sec:reason-measurable}
&
simple empirical laws
&
\celllist{
neural scaling laws\\
edge of stability\\
neural feature ansatz
}
&
\celllist{
the laws of Kepler, Snell, Boyle, Hooke, Newton, Faraday, Ohm, Poiseuille, Planck, Hubble, etc.
}
\\[3em]

\ref{sec:reason-hps}
&
\celllist{
study of system\\
parameters
}
&
\celllist{
step size as sharpness\\
regularization,\\
$\mu$P and width-scaling
}
&
\celllist{
scaling analysis,\\
nondimensionalization,\\
chaotic vs. ordered regimes
}
\\[3em]

\ref{sec:reason-universal-behavior}
&
universal phenomena
&
\celllist{
common inductive biases and representations across models
}
&
\celllist{
critical phenomena,\\
renormalization group flow
}
\\[2em]

\bottomrule
\end{tabular}

\caption{
\textbf{Useful tools and ideas in the emerging science of deep learning closely resemble important tools and ideas from physics, particularly classical mechanics, continuum mechanics, statistical mechanics, and quantum mechanics.}
Extrapolating, this suggests that there will be a \textit{mechanics of learning} which offers a unifying first-principles theory of the training process, hidden representations, final weights, and test-time performance of neural networks.
}
\label{tab:lm_physics_comparison}
\end{table}
